% Tez ana dosyası. Derleme: pdfLaTeX + BibTeX (Overleaf varsayılanı) ya da latexmk -pdf main.tex.
% Figürler ../results/ klasöründen okunur; Overleaf için make_overleaf_zip.py figürleri figures/ altına toplar.
% Taslak 2 (kısa, Nejakar 2025 iskeleti), 2026-11-22. Uzun ilk taslak: arsiv_uzun/ (detay gerektikçe oradan). Sayılar tez notlarının "FINAL REGENERATION 2026-11-22" bölümünden.
\documentclass[11pt,a4paper,oneside]{report}

\usepackage[utf8]{inputenc}
\usepackage[T1]{fontenc}
\usepackage{lmodern}
\usepackage[english]{babel}
\usepackage[a4paper,margin=2.5cm]{geometry}
\usepackage{setspace}\onehalfspacing
\usepackage{amsmath,amssymb}
\usepackage{siunitx}
\sisetup{detect-all, per-mode=symbol}
\usepackage{graphicx}
\graphicspath{{figures/}{../results/}}
\usepackage{booktabs,tabularx,array}
\usepackage{caption}\captionsetup{font=small,labelfont=bf}
\usepackage[numbers,sort&compress]{natbib}
\usepackage{xcolor}
\usepackage[hidelinks]{hyperref}
\usepackage{cleveref}
\usepackage{enumitem}\setlist{noitemsep,topsep=3pt}

% taslak notu: son halde \todo{} satırları silinir
\newcommand{\todo}[1]{\textcolor{orange}{[TODO: #1]}}
% figür kısayolu, argümanlar: dosya adı, genişlik (\linewidth oranı), başlık, etiket
\newcommand{\fig}[4]{\begin{figure}[htbp]\centering\includegraphics[width=#2\linewidth]{#1}\caption{#3}\label{#4}\end{figure}}

\title{Digital-Twin-Based Battery Management System\\with Cloud Monitoring and SOC Estimation}
\author{Berkan Talha Yanbaş}
\date{\todo{month year}}

\begin{document}

\begin{titlepage}
  \centering
  {\large \todo{University}\\ \todo{Faculty / Department}\par}
  \vspace{3cm}
  {\LARGE\bfseries Digital-Twin-Based Battery Management System with Cloud Monitoring and SOC Estimation\par}
  \vspace{2cm}
  {\large Berkan Talha Yanbaş\par}
  \vspace{1cm}
  {\large Bachelor's Thesis\par}
  \vfill
  {\large Supervisor: \todo{name}\par}
  \vspace{1cm}
  {\large \todo{City, month year}\par}
\end{titlepage}

\pagenumbering{roman}
\chapter*{Abstract}
\addcontentsline{toc}{chapter}{Abstract}

This thesis presents a digital-twin-based battery management system (BMS) with cloud monitoring for a 19S12P electric-vehicle pack of LG INR21700 M50-family cells. A digital twin of the pack, parameterised with published M50T measurements, runs on an Arduino Uno and sends its measurements over a real CAN bus to the BMS, which runs on an STM32G431 microcontroller. The BMS estimates the state of charge (SOC) with an extended Kalman filter that also tracks the current-sensor offset and the cell resistance, estimates the core temperature, detects unexplained heating, and limits the allowed current in four graded states. A PC gateway forwards the telemetry over HTTPS to a web dashboard.

The SOC error is 0.277\,\% RMSE on the reference drive. In 1000 randomised closed-loop runs on the microcontroller every fault was handled as specified, an external heat source was detected after a median of 8.0\,s, and the cloud link delivered all 18\,046 rows of a 30-minute test with a median round-trip time of 102\,ms.

\textbf{Keywords:} battery management system, digital twin, state of charge, extended Kalman filter, CAN, cloud monitoring.

\chapter*{Özet}
\addcontentsline{toc}{chapter}{Özet}

Bu tezde, 19S12P LG INR21700 M50 ailesi hücrelerden oluşan bir elektrikli araç paketi için dijital ikiz tabanlı, bulut izlemeli bir batarya yönetim sistemi (BMS) sunulmaktadır. Yayımlanmış M50T ölçümleriyle parametrelenen paket ikizi bir Arduino Uno üzerinde çalışır ve ölçümlerini gerçek bir CAN hattı üzerinden STM32G431 mikrodenetleyicideki BMS'e gönderir. BMS, akım sensörü ofsetini ve hücre direncini de izleyen genişletilmiş Kalman filtresiyle şarj durumunu (SOC) kestirir, çekirdek sıcaklığını tahmin eder, açıklanamayan ısınmayı fark eder ve izin verilen akımı dört kademeli durumla sınırlar. Bir PC ağ geçidi telemetriyi HTTPS üzerinden web paneline iletir.

SOC hatası referans sürüşte 0,277\,\% RMSE'dir. Mikrodenetleyicide koşturulan 1000 rastgele kapalı çevrim koşusunda her arıza tanımlandığı gibi ele alınmış, dış ısı kaynağı ortalama 8,0\,s'de fark edilmiş ve bulut bağlantısı 30 dakikalık testte 18\,046 satırın tamamını 102\,ms ortanca gidiş-dönüş süresiyle iletmiştir.

\textbf{Anahtar kelimeler:} batarya yönetim sistemi, dijital ikiz, şarj durumu, genişletilmiş Kalman filtresi, CAN, bulut izleme.

\tableofcontents
\listoffigures
\cleardoublepage
\pagenumbering{arabic}

\chapter{Introduction}\label{ch:intro}

Electric vehicles depend on lithium-ion packs that are safe only inside a narrow window of voltage, current and temperature \citep{feng2018}. The battery management system (BMS) keeps the pack inside this window: it measures cell voltages, current and temperature, estimates hidden quantities such as the state of charge (SOC) and the core temperature, and decides how much current the vehicle may use.

Recent work combines the BMS with a \emph{digital twin}, a virtual replica of the battery that is kept in sync with the real one, and with cloud services for remote monitoring \citep{li2020,nejakar2025}. Nejakar et al.\ built such a system with an ESP32, low-cost sensors, a Firebase database and a web dashboard, and estimated the SOC with a hybrid Kalman filter and recurrent neural network \citep{nejakar2025}. Their system monitors and raises threshold alerts, but it does not control the battery, and it is tested on a single 12\,V battery.

This thesis takes the same three building blocks (digital twin, SOC estimation, cloud monitoring) and adds the part a vehicle needs: a BMS controller that limits the current and reacts safely to faults. The digital twin is used in a second role as well: it replaces the physical pack on the test bench, so that faults which are too dangerous for real cells can be tested repeatably.

\section*{Contributions}
\begin{enumerate}
  \item A BMS controller for a 19S12P pack on an STM32 microcontroller: SOC estimation with an extended Kalman filter (EKF), core-temperature estimation, detection of unexplained heating and a graded current limit (NORMAL, DERATE, LIMP-HOME, SAFE-STOP).
  \item A digital twin of the pack on a separate microcontroller, fitted to published measurements of the M50T cell \citep{kirkaldy2024} and connected only through CAN.
  \item Cloud monitoring through an HTTPS gateway and a web dashboard, with measured link quality.
  \item Verification from simulation to the hardware bench, including 1000 randomised runs on the microcontroller.
\end{enumerate}

\Cref{ch:background} gives the background, \cref{ch:method} the methodology, \cref{ch:hardware} the hardware, \cref{ch:results} the results and \cref{ch:conclusion} the conclusion.

\chapter{Background}\label{ch:background}

\paragraph{SOC estimation.} Coulomb counting integrates the current; it is cheap but drifts with sensor offset and noise. Reading the SOC from the open-circuit voltage (OCV) works only after the cell has rested. Model-based estimators combine both: an extended Kalman filter (EKF) predicts the cell state with an equivalent circuit model and corrects it with the measured voltage \citep{plett2004ekf}. The model used here is a resistance $R_0$ in series with two resistor--capacitor branches; all resistances depend on temperature through the Arrhenius law \citep{plett2015}. A sensor offset can be added to the filter state, but it is observable only where the OCV curve is curved \citep{zhao2017}. Data-driven estimators such as recurrent neural networks can learn the non-linear behaviour from history \citep{nejakar2025}.

\paragraph{Thermal behaviour.} Heat generated in the cell is $q = i\,(\mathrm{OCV} - v)$ \citep{bernardi1985}. A cylindrical cell is modelled by two thermal nodes, core and surface \citep{lin2014}. Only the surface can be measured, so the core temperature has to be estimated.

\paragraph{Digital twins and cloud BMS.} A digital twin mirrors the battery's state in software. Li et al.\ ran SOC and state-of-health estimation in the cloud \citep{li2020}; Nejakar et al.\ synchronised an ESP32 monitor with a cloud database and a dashboard \citep{nejakar2025}. Connecting a BMS to the cloud also opens security risks \citep{naseri2023}, so safety decisions should stay on the vehicle.

\paragraph{Hardware-in-the-loop testing.} In hardware-in-the-loop (HIL) testing a real controller runs against a simulated plant. Laboratory BMS rigs use real-time simulators and cell-voltage emulators \citep{barreras2016,bui2019}. When only the control strategy is tested, it is sufficient to simulate the cell modules over CAN \citep{haupt2013}, which is the approach taken here.

\chapter{Methodology}\label{ch:method}

The system has four layers (\cref{fig:arch}): the \emph{physical layer} is the digital twin of the pack, the \emph{edge layer} is the BMS controller, the \emph{cloud layer} is the gateway and server, and the \emph{user interface} is the web dashboard. The methodology is described in six blocks, following the structure of \citet{nejakar2025}.

\fig{system_block_diagram}{0.95}{System architecture. \todo{redraw as four layers (twin / BMS / cloud / dashboard), English labels}}{fig:arch}

\section{Digital twin of the pack}
The twin simulates a 19S12P pack (58.2\,Ah, about 4\,kWh): 19 series groups, each with its own SOC and a two-RC electrical model, a two-node thermal model and a simple vehicle model that turns a drive cycle into a current demand. Its parameters come from public measurements of 40 LG M50T cells \citep{kirkaldy2024}: the OCV curve, the resistances, their temperature dependence and the resistance growth with ageing. The twin adds realistic sensor noise and clips its current to the limit sent by the BMS, which closes the loop. It runs at 10\,Hz on an Arduino Uno.

\section{Real-time data acquisition}
Every 100\,ms the twin sends the pack current, pack voltage, surface and ambient temperature and the 19 group voltages over CAN at 500\,kbit/s. Each frame carries a counter and a CRC-8 checksum, so lost, repeated or corrupted frames are detected; if no valid measurement arrives for 500\,ms the BMS treats the measurement line as lost.

\section{SOC estimation}
The SOC is estimated with an EKF whose state is
\begin{equation}
  x = [\,z,\ v_1,\ v_2,\ b_I,\ k_R\,]^\top,
\end{equation}
the SOC $z$, the two RC voltages, the current-sensor offset $b_I$ and a resistance scale $k_R$. The offset state removes the drift of Coulomb counting; the resistance state keeps the filter accurate as the cell ages. At start-up the contactor stays open for one second, so the measured current is pure offset; the filter starts from this value and an offset above 3\,A is reported as a sensor fault.

As a machine-learning component, a small decision tree trained on the same ageing data estimates the state of health (SOH) from usage and resistance and runs on the microcontroller. \todo{decide how much ML to show}

\section{Cloud integration}
The BMS writes one telemetry line per step to its serial port (SOC, SOH, temperatures, allowed current, state) and an event line whenever its state changes. A Python gateway on a PC groups the lines into 1\,s batches and sends them over HTTPS to a PHP server with an SQLite database. The gateway authenticates with an API key and numbers each batch, so the server rejects repeated batches.

\section{Web dashboard}
The dashboard shows the live state, SOC, temperatures, allowed and measured current and the reasons for any current limit. A second view shows the stress criteria and the reaction times of the bench fault tests.

\section{Decision and alert mechanism}
Instead of fixed alarm thresholds, the BMS limits the current in four graded states:
\begin{description}
  \item[NORMAL] full datasheet current;
  \item[DERATE] reduced current because of high temperature, a cooling problem, a fast temperature rise or low voltage headroom;
  \item[LIMP-HOME] 25\,\% current because a sensor failed or unexplained heating was detected;
  \item[SAFE-STOP] zero current and the contactor opened: core above 60\,°C, voltage outside its limits, lost measurements, or a software fault.
\end{description}
The derating criteria are combined into one stress index $\Psi = \min\big(1, \Psi_\mathrm{max} + 0.25\,(\sum\Psi_i - \Psi_\mathrm{max})\big)$ and the allowed current is $I_\mathrm{datasheet}(1-\Psi)$.

Two estimators feed this logic. A second Kalman filter estimates the core temperature and the cooling resistance, which reveals a failing cooling system. A thermal anomaly detector checks whether the measured surface temperature rises faster than the known heat could explain even without cooling; if so, an unknown heat source is present and the BMS goes to LIMP-HOME.

\chapter{Hardware setup}\label{ch:hardware}

\Cref{tab:hw} lists the bench. The BMS and the twin run on separate boards and share only the CAN bus, so the BMS sees the same kind of data it would get from a real pack.

\begin{table}[htbp]\centering\small
\caption{Hardware of the test bench.}\label{tab:hw}
\begin{tabularx}{\linewidth}{@{}>{\raggedright\arraybackslash}p{0.34\linewidth}>{\raggedright\arraybackslash}X@{}}\toprule
Component & Role \\\midrule
STM32 NUCLEO-G431RB (Cortex-M4F, 170\,MHz) & BMS controller; floating-point unit, CAN controller, watchdog \citep{rm0440} \\
Arduino Uno + MCP2515 & Digital twin of the pack; CAN controller over SPI \\
2 $\times$ TJA1050 & CAN transceivers, 500\,kbit/s \citep{tja1050} \\
LM35 & Ambient temperature for the twin \\
Push buttons, potentiometer & Fault scenarios and their severity \\
I$^2$C LCD & Local display of state, SOC and allowed current \\
LEDs via ULN2003, relay & ``Headlights'' showing the allowed current; relay as main contactor \\
PC & Gateway to the cloud server \\\bottomrule
\end{tabularx}
\end{table}

\todo{bench photo}

The BMS firmware is written in C at register level without a vendor library. One control step takes at most 92\,\textmu s, less than 0.1\,\% of the 100\,ms period. The relay is held closed only while the BMS actively drives it, so a crash or power loss disconnects the pack. A watchdog resets a hung controller, and a checksum of the firmware is verified at every start.

\chapter{Results and discussion}\label{ch:results}

\section{Model accuracy}
The twin is only as good as its cell model. The OCV curve was rebuilt from the GITT tests of the M50T data set and validated on cells that were not used to build it: the voltage error fell from 64.7\,mV with a literature curve for the sister cell \citep{chen2020} to 17.0\,mV (\cref{fig:ocv}). The temperature dependence of the resistances and the resistance growth with ageing ($R/R_\mathrm{new} = 1 + 1.31\,(1-\mathrm{SOH})$, 40 cells) were also fitted to the measurements.

\fig{ocv_dogrulama}{0.85}{OCV validation on cells not used for fitting.}{fig:ocv}

\section{SOC estimation accuracy}
\Cref{tab:soc} compares SOC estimators on the same reference drive. The EKF with temperature-dependent resistance and offset state has the lowest error.

\begin{table}[htbp]\centering\small
\caption{SOC estimation error on the reference drive (after 60\,s).}\label{tab:soc}
\begin{tabular}{@{}lrr@{}}\toprule
Method & RMSE (\%) & Max.\ error (\%) \\\midrule
OCV lookup & 17.0 & 36.2 \\
Coulomb counting & 0.86 & 1.14 \\
EKF, fixed resistance & 0.72 & 1.09 \\
\textbf{EKF, $R(T)$ + offset (proposed)} & \textbf{0.28} & \textbf{0.52} \\\bottomrule
\end{tabular}
\end{table}

\fig{estimator_comparison_normal}{0.85}{SOC error of the estimators.}{fig:soc}

The start-up offset check detects a current-sensor offset of 4\,A or more within 0.9\,s; without it the filter needed more than 1000\,s or did not detect it at all.

\section{Fault detection and safety}
The twin and the BMS were compiled together and run on the STM32 in 1000 randomised runs of 600\,s each, with random temperature, SOC, driving, cell spread and one fault per run. Every run passed its verdict, and a PC build of the same code gave the same verdict in all 1000 runs (\cref{tab:mc}).

\begin{table}[htbp]\centering\small
\caption{On-target Monte Carlo test, 1000 runs.}\label{tab:mc}
\begin{tabular}{@{}lrl@{}}\toprule
Fault & Runs & Reaction \\\midrule
None & 178 & no false alarm \\
External heat & 151 & detected after 8.0\,s (median) \\
Temperature sensor open & 151 & detected in the same step \\
Current-sensor offset & 202 & detected after 0.9\,s \\
Cooling failure & 199 & detected in 48\,\% (median 138\,s), always safe \\
Aged pack & 119 & resistance growth recognised \\\bottomrule
\end{tabular}
\end{table}

\fig{mc_overview}{0.9}{Monte Carlo results per fault type.}{fig:mc}

On the bench, the hardware tests found three design flaws that simulation had missed: a cooling alarm that switched on and off repeatedly, a false cooling alarm after restarting the BMS on a warm pack, and a core-temperature estimate that lagged far behind the truth under external heat. All three were fixed; with the fix, the safe stop now comes at a true core temperature of 60\,°C instead of 66\,°C.

\section{Cloud synchronisation}
Over 30 minutes the gateway delivered all 18\,046 telemetry rows to the server without loss. The median upload round-trip time was 102\,ms (95th percentile 134\,ms); the end-to-end delay of about 1\,s is mostly the batching interval (\cref{fig:link}).

\fig{link_quality_20261121}{0.85}{Cloud link over 30 minutes.}{fig:link}

\section{Dashboard}
\Cref{fig:dash} shows the dashboard during a bench test. It shows the state and the reason for any current limit, so a user can see why the battery is being protected.

\fig{web_dashboard_live}{0.85}{Web dashboard.}{fig:dash}

\section{Comparison with the base paper}
\Cref{tab:cmp} compares this work with \citet{nejakar2025}. The two SOC errors are not directly comparable: theirs is a mean absolute error on a real 12\,V battery, ours a root-mean-square error against the twin's true SOC.

\begin{table}[htbp]\centering\small
\caption{Comparison with the base paper.}\label{tab:cmp}
\begin{tabularx}{\linewidth}{@{}l>{\raggedright\arraybackslash}X>{\raggedright\arraybackslash}X@{}}\toprule
 & \citet{nejakar2025} & This work \\\midrule
Battery & real 12\,V battery & digital twin of a 19S12P pack (M50T data) \\
Controller & ESP32 (monitoring) & STM32 BMS (monitoring + control) \\
Link & Wi-Fi to Firebase & CAN to BMS, HTTPS to own server \\
SOC method & Kalman + RNN, MAE 2.6\,\% & EKF with offset and resistance states, RMSE 0.28\,\% \\
Alerts & thresholds (SOC $<$ 20\,\%, T $>$ 50\,°C) & four graded states, stress index, anomaly detection \\
Cloud test & 48\,h, no loss, latency $<$ 1\,s & 30\,min, no loss, RTT 102\,ms \\
Validation & real use & 1000 on-target runs + HIL bench \\\bottomrule
\end{tabularx}
\end{table}

\section{Limitations}
The twin sends its measurements over CAN, so the analog measurement hardware of a real BMS is not tested. The BMS does not estimate capacity, so on an aged pack the SOC error rises to about 5\,\%. Results come from a simulated pack; a test on real cells remains future work.

\chapter{Conclusion}\label{ch:conclusion}

This thesis built a digital-twin-based BMS with cloud monitoring. A twin of a 19S12P pack, fitted to measured M50T data, feeds a BMS controller on an STM32 over CAN, and a gateway forwards the telemetry to a web dashboard. The BMS estimates the SOC with 0.28\,\% RMSE, detects unexplained heating within seconds and limits the current in four graded states. It passed 1000 randomised runs on the microcontroller, and the cloud link delivered all telemetry without loss.

Compared with cloud monitoring systems such as \citet{nejakar2025}, the main step forward is control: the BMS does not only report a problem, it protects the battery, and the digital twin makes it possible to test this safely.

\paragraph{Future work.} Testing with real cells and a cell-monitoring board; capacity estimation; a learned SOC component combined with the EKF, as in the base paper; and a mobile interface.


\bibliographystyle{unsrtnat}
\bibliography{references}


\end{document}
